\documentclass[11pt]{article}

\usepackage[margin=1in]{geometry}
\usepackage{amsmath,amssymb,amsthm,mathtools}
\usepackage{booktabs}
\usepackage{array}
\usepackage{enumitem}
\usepackage{xcolor}
\usepackage[most]{tcolorbox}
\usepackage{hyperref}
\hypersetup{colorlinks=true,linkcolor=teal!55!black,urlcolor=teal!55!black,citecolor=teal!55!black}

\definecolor{navy}{HTML}{13233F}
\definecolor{teal}{HTML}{1C7293}
\definecolor{amber}{HTML}{C77700}
\definecolor{plainbg}{HTML}{F4F1E8}
\definecolor{plainfr}{HTML}{B8923A}
\definecolor{boxbg}{HTML}{EEF3F8}
\definecolor{ambbg}{HTML}{FFF6E8}
\definecolor{okbg}{HTML}{EAF5EE}
\definecolor{okfr}{HTML}{2A9D5C}

% Plain-English box
\newtcolorbox{plain}{colback=plainbg,colframe=plainfr,boxrule=1pt,arc=2pt,
  left=8pt,right=8pt,top=6pt,bottom=6pt,
  fonttitle=\bfseries,title=In plain English}
\newtcolorbox{keybox}{colback=boxbg,colframe=teal,boxrule=1pt,arc=2pt,left=8pt,right=8pt,top=6pt,bottom=6pt}
\newtcolorbox{warnbox}{colback=ambbg,colframe=amber,boxrule=1pt,arc=2pt,left=8pt,right=8pt,top=6pt,bottom=6pt}
\newtcolorbox{okbox}{colback=okbg,colframe=okfr,boxrule=1pt,arc=2pt,left=8pt,right=8pt,top=6pt,bottom=6pt}

\theoremstyle{definition}
\newtheorem{assumption}{Assumption}
\theoremstyle{plain}
\newtheorem{lemma}{Lemma}
\newtheorem{theorem}{Theorem}

\newcommand{\Fxeb}{F_{\mathrm{XEB}}}
\newcommand{\Hmin}{H_{\min}}
\newcommand{\epss}{\varepsilon_{s}}
\newcommand{\epssou}{\varepsilon_{\mathrm{sou}}}
\newcommand{\Qmin}{Q_{\min}}
\newcommand{\E}{\mathbb{E}}

\title{\vspace{-1.5em}\textbf{Certified Quantum Randomness}\\[0.3em]
\large Protocol, Mathematics, and Status --- a consolidated reference}
\author{RCS-VRF Subnet --- working notes}
\date{\today}

\begin{document}
\maketitle
\vspace{-1em}

\begin{keybox}
\textbf{What this document is.} A single ordered reference pulling together the
intuition \emph{and} the mathematics behind the RCS-VRF certified-randomness protocol:
why it exists, what it must achieve, how the XEB verification works (with a fully
worked numerical example), the certified min-entropy bound (anchored to Liu \emph{et al.},
Nature \textbf{640}:343, 2025; arXiv:2503.20498), and an honest accounting of what is
proven versus still open. Plain-English boxes summarise each technical section.
\end{keybox}

\tableofcontents
\bigskip

% =====================================================================
\part*{Part I --- The big picture}
\addcontentsline{toc}{section}{Part I --- The big picture}
% =====================================================================

\section{What the protocol is for}
The goal is to produce random bits that come with a \emph{proof} they are genuinely
unpredictable and genuinely quantum in origin --- verifiable by anyone with a laptop,
without trusting the hardware provider. A client generates pseudo-random ``challenge''
circuits from a small seed, sends them to an untrusted quantum server, and asks for one
measured bitstring per circuit. The server's results are then checked, and if they pass,
distilled into certified-uniform randomness.

\begin{plain}
We sell ``randomness you don't have to trust us for.'' Ordinary random-number generators
say ``trust me, it's random.'' This protocol instead hands you a number \emph{and} a
mathematical certificate that no classical computer could have produced it on demand --- so
a casino, lottery, or blockchain can verify fairness without trusting the operator.
\end{plain}

\section{The four required properties}
A batch of output must satisfy four properties before it can be sold as certified quantum
randomness:
\begin{enumerate}[leftmargin=1.6em,itemsep=2pt]
\item \textbf{Quantum in origin} --- the bits really came from quantum hardware, not a
classical RNG. (Established by the XEB score plus the timing deadline.)
\item \textbf{Genuinely random} --- truly unpredictable and uniform, not merely
``quantum-flavoured.'' (Min-entropy bound $+$ randomness extraction.)
\item \textbf{Classically verifiable} --- anyone can check properties 1--2 with only a
laptop and the recorded log. (Recompute the score, replay the audit log.)
\item \textbf{On-demand} --- fresh randomness whenever requested, quickly. (The
unpredictable seed drops, the answer returns inside a deadline.)
\end{enumerate}

\begin{plain}
Where did it come from (quantum), is it actually random (provably yes), can outsiders
check (yes, classically), can you get it when you ask (yes, fast). Properties 1 and 4 are
about the \emph{source}; properties 2 and 3 are about the \emph{guarantee}. They prop each
other up: on-demand is \emph{how} we get quantum-in-origin (the deadline rules out
faking), and verifiability is \emph{how} anyone trusts the rest.
\end{plain}

\section{Who is the ``attacker,'' and why we model one}
The attacker is not an outside hacker --- it is a \emph{miner}: a participant paid to run
the quantum sampling who would rather collect the reward without doing the expensive
quantum work (by faking samples classically), or who wants to \emph{bias} the output to
win a bet. We assume the worst-case: the most capable, most motivated such participant.

\begin{plain}
Defending against the attacker \emph{is} the product. Strip it out and you've built an
expensive RNG that says ``trust me.'' The value is precisely that the output holds up
against a participant with money on the line and every reason to cheat. We assume the
worst case because security that only works against lazy attackers is luck, not security.
\end{plain}

\section{Proving ``quantum'' by ruling out ``classical''}
We cannot directly inspect qubits. Instead we rule out every \emph{feasible classical}
way the samples could have been produced; what remains is genuine quantum sampling. The
list of ``classical methods to rule out'' is identical to the ``attacker playbook,'' so
one argument does both jobs.

\begin{warnbox}
The claim is ``not \emph{feasibly} classical,'' not ``not classical in principle.'' A
classical computer \emph{can} reproduce the samples given enough time; we only prove it
cannot do so \emph{here, now, on demand}, within the deadline. The proof is therefore only
as strong as our list of classical strategies being complete --- which is why the hardness
of spoofing XEB is an active research assumption, not a theorem.
\end{warnbox}

\section{The timing asymmetry: ``a clock, not a computer''}
The attacker and the verifier face the \emph{same} hard calculation, but under opposite
time pressure. The attacker must produce passing samples \emph{fast and on demand}; the
verifier only has to \emph{recognise} them \emph{slowly and offline, after the fact}.

\begin{plain}
The quantum machine answers in seconds (it just measures). A classical faker must simulate
the circuit, which takes hours--days, so it misses the deadline. The verifier checks the
recorded log later with no clock running, so it can take all the time it needs. Two legs
exist purely to police this clock: the VRF seed stops the attacker starting early, the
timing check enforces the deadline. The security is the \emph{gap in deadlines}, not a gap
in horsepower.
\end{plain}

% =====================================================================
\part*{Part II --- How verification works (the XEB machinery)}
\addcontentsline{toc}{section}{Part II --- How verification works}
% =====================================================================

\section{Two roles: prediction vs.\ measurement}
The same distribution --- the circuit's ideal output probabilities
$p_C(x)=|\langle x|C|0\rangle|^2$ --- is used by both parties, in opposite senses.
\begin{itemize}[leftmargin=1.6em,itemsep=1pt]
\item \textbf{Step 0 (verifier, prediction).} Computes the probabilities by classically
simulating the circuit. Knows the whole ``map,'' but cannot sample cheaply. Deterministic;
contains \emph{no} randomness. Exponentially expensive.
\item \textbf{Step 1 (server, measurement).} Runs the circuit and measures, collapsing the
state to one bitstring with probability exactly $p_C(x)$ (Born rule). Effortless for a
quantum device, but yields one draw and no knowledge of the odds. This is the \emph{only}
step where randomness enters.
\end{itemize}

\begin{plain}
Step 0 is the verifier drawing a perfect \emph{weather map} --- every region's rain
probability, computed exactly but at great cost. Step 1 is the server reporting whether
it's \emph{actually raining right now} --- one real observation, instant, but unpredictable
even to someone holding the perfect map. The protocol works because the map is expensive
to draw (a faker can't quickly forge one) and the actual weather is genuinely
unpredictable (no map pre-determines it). Even the verifier, who knows the entire
distribution, cannot predict the outcome of a single measurement --- that residual
unpredictability is the certified randomness.
\end{plain}

\subsection{Why sampling is easy but prediction is hard}
Sampling (Step 1) is free for a quantum computer --- it is what the hardware physically
does. Predicting (Step 0) means computing amplitudes, which is exponentially costly. A
classical faker is asked to do the easy-for-quantum thing (sample) and the only classical
route is to first do the hard-for-classical thing (predict). That asymmetry is the resource
the protocol monetises.

\section{How a circuit is simulated (Step 0 in detail)}
Two clarifications:
\begin{itemize}[leftmargin=1.6em,itemsep=2pt]
\item \textbf{The gates are not random at simulation time.} The circuit is pseudo-random
but generated \emph{deterministically from the seed} (single-qubit $SU(2)$ gates from a
seeded PRG; two-qubit topology fixed by an edge-colouring). Once the seed is fixed the
circuit is one concrete, fully-determined object that verifier and server both reconstruct.
\item \textbf{Simulation computes one exact amplitude, run once --- not many shots.} The
verifier needs $p_C(x_i)$ only for the specific strings the server returned. For each it
computes the single number $\langle x_i|C|0\rangle$ \emph{exactly, once}. No statistics, no
repetition.
\end{itemize}
Computing $\langle x|C|0\rangle$ for $n=56$ cannot use the naive $2^n$ state vector
($2^{56}\approx7\times10^{16}$ amplitudes). Instead one uses \textbf{tensor-network
contraction}: represent each gate as a small tensor, fix input $|0\rangle$ and output $x$,
and contract the network down to a single scalar without ever materialising the full state.
Index slicing breaks the contraction into GPU-sized pieces.

\begin{keybox}
The per-circuit simulation cost is $B\approx 90\times10^{18}$ FLOPs --- about $100$ seconds
on the full Frontier supercomputer for \emph{one} amplitude. Crucially, $B$ is both the
verifier's cost (paid offline, no clock) \emph{and} the cost a classical faker would pay to
predict the distribution (which it must pay \emph{within the deadline}). That shared cost,
under opposite time pressure, is the security argument made quantitative.
\end{keybox}

\section{The cross-entropy benchmark $\Fxeb$}
\label{sec:xebdef}
The verification statistic is the linear cross-entropy benchmark, over a random test subset
$V$ of size $m$:
\begin{equation}\label{eq:xebdef}
\Fxeb \;=\; \frac{N}{m}\sum_{i\in V} p_{C_i}(x_i)\;-\;1, \qquad N=2^n .
\end{equation}
Up to constants it is the \emph{average ideal-probability of the strings the server
returned}, rescaled so that blind guessing scores $0$ and perfect quantum sampling scores
$1$.

\begin{plain}
$\Fxeb$ asks: ``do the strings the server sent land on the circuit's \emph{likely}
outcomes more than chance would?'' An honest quantum machine returns high-probability
strings (the Born rule biases it that way), so the average is high. A blind guesser
scatters everywhere and averages out to chance. The score is built so chance $=0$ and
perfect $=1$.
\end{plain}

\subsection{A fully worked example ($n=2$)}
Take $N=2^2=4$, uniform baseline $1/N=0.25$. Suppose the verifier's simulation (Step 0)
gives
\[
p_C(00)=0.50,\quad p_C(01)=0.30,\quad p_C(10)=0.15,\quad p_C(11)=0.05 .
\]
The server returns $m=5$ samples (Step 1): $\;00,\,01,\,00,\,11,\,01$. Then:
\begin{center}
\renewcommand{\arraystretch}{1.15}
\begin{tabular}{r l l}
\toprule
Step & Operation & Result \\
\midrule
2 & look up each returned string's $p_C$ & $0.50,0.30,0.50,0.05,0.30$ \\
3 & sum & $\sum p_C(x_i)=1.65$ \\
4 & average ($\div m$) & $1.65/5 = 0.33$ \\
5 & rescale ($\times N$) so chance $=1$ & $4\times0.33 = 1.32$ \\
6 & subtract $1$ so chance $=0$ & $\Fxeb = 0.32$ \\
\bottomrule
\end{tabular}
\end{center}
Each step has one job: Step 2 connects samples to the circuit, Step 4 averages out
per-sample noise, Step 5 sets the unit, Step 6 sets the zero. Sanity checks on the same
circuit: blind guessing (all four strings equally) averages $p_C\to0.25$, giving
$\Fxeb=4(0.25)-1=0$; returning only the top string $00$ gives $4(0.50)-1=1$.

\begin{plain}
Strings come in $\to$ look up each one's ideal probability $\to$ average them $\to$ rescale
so chance $=1$ $\to$ subtract so chance $=0$. The arithmetic is trivial; the expensive part
is Step 2's lookups, because computing each $p_C$ means simulating the circuit.
\end{plain}

\section{The mathematics of the reference cases}
Everything follows from the \textbf{Porter--Thomas} law for a random circuit's output
probabilities,
\begin{equation}\label{eq:pt}
f(p) = N e^{-Np}, \qquad \E[p]=\frac1N .
\end{equation}

\paragraph{$\Fxeb=0$ (uniform / spoof).} A string chosen independently of the circuit lands
on the plain PT mean:
\[
\E_{x\sim U}[p_C(x)] = \int_0^\infty p\,f(p)\,dp = \frac1N
\;\Longrightarrow\;
\Fxeb=\frac{N}{m}\,m\,\tfrac1N - 1 = 0 .
\]

\paragraph{$\Fxeb=1$ (perfect quantum).} Quantum sampling returns $x$ with probability
$p_C(x)$, \emph{size-biasing} the density by an extra factor $p$:
\[
\mathrm{PDF}_{x\sim Q}[p]=N\,f(p)\,p = N^2 p\,e^{-Np},\qquad
\E_{x\sim Q}[p_C(x)] = \int_0^\infty p\,(N^2 p e^{-Np})\,dp = N^2\frac{2}{N^3}=\frac2N .
\]
The sampled probability is \emph{twice} the uniform mean --- that factor of $2$ is the
quantum signal:
\[
\Fxeb=\frac{N}{m}\,m\,\tfrac2N - 1 = 1 .
\]

\paragraph{$0<\Fxeb=\phi$ (noisy real machine).} A device of fidelity $\phi$ is modelled
(depolarising) as $\rho_\phi=\phi\,|\psi\rangle\langle\psi|+(1-\phi)\,I/N$, whose sampling
PDF is the convex mixture
\[
\mathrm{PDF}_{x\sim Q_\phi}[p]=\phi\,\mathrm{PDF}_Q[p]+(1-\phi)\,\mathrm{PDF}_U[p],
\qquad
\E[p_C(x)]=\phi\tfrac2N+(1-\phi)\tfrac1N=\frac{1+\phi}{N}.
\]
Hence the headline relationship
\begin{equation}\label{eq:fxebphi}
\boxed{\;\Fxeb = (1+\phi)-1 = \phi\;}
\end{equation}
The score \emph{equals the fidelity}: Liu \emph{et al.}\ had $\phi\approx0.3$ and measured
$\Fxeb=0.32$.

\paragraph{$\Fxeb>1$ (cherry-picking --- why the band is two-sided).} A cheater that computes
$p_C$ and returns only the \emph{highest}-probability strings reports the expected maximum,
$\E[\max_x p_C(x)]\approx \tfrac{\ln N}{N}=\tfrac{n\ln2}{N}$, giving
$\Fxeb\approx n\ln2-1>1$. No physical noisy device can exceed $\sim1$ (noise only pulls
down). Hence the acceptance rule is two-sided,
\[
\chi_{\text{low}} \le \Fxeb \le \chi_{\text{high}},
\]
the low edge killing uniform spoofers, the high edge killing cherry-pickers.

\begin{warnbox}
$\Fxeb\approx\phi$ holds only in the \emph{low-noise} regime ($\epsilon n<\ln3\approx1.1$;
the experiment sits at $\epsilon n\approx0.13$). Past that threshold score and fidelity
decouple --- one can get high $\Fxeb$ at low fidelity. So ``score $=$ fidelity'' is a
regime-dependent truth, not an identity.
\end{warnbox}

\subsection{The full distributions (used by the security proof)}
The means above are the centres; the proof needs the CDFs. Because $\sum_i p_{C_i}(x_i)$ is
a sum of exponential-type variables, the XEB score follows \textbf{Erlang}/regularised
incomplete-gamma ($\widetilde\Gamma$) laws:
\[
\Pr[\Fxeb\le\chi] =
\begin{cases}
\widetilde\Gamma(m,\,m(\chi+1)) & \text{uniform ($0$ quantum samples),}\\[2pt]
\widetilde\Gamma(2m,\,m(\chi+1)) & \text{perfect quantum,}\\[2pt]
\widetilde\Gamma(m+l,\,m(\chi+1)) & \text{mixture of $l$ quantum }+\,(m-l)\text{ uniform.}
\end{cases}
\]
The mixture form $\widetilde\Gamma(m+l,\dots)$ is the workhorse of the $\Qmin$ calculation
(Part III).

% =====================================================================
\part*{Part III --- The certified entropy (Leg 5)}
\addcontentsline{toc}{section}{Part III --- The certified entropy}
% =====================================================================

\section{Why min-entropy, and extraction}
Passing XEB shows the samples are quantum, but they are not yet \emph{uniform} --- a real
device's bits are biased and partly predictable. The relevant currency is the
\textbf{smooth min-entropy} $\Hmin^{\epss}$, which measures how many truly-uniform bits can
be squeezed out. A randomness extractor (here a Toeplitz hash) then compresses the raw bits
to that many uniform bits. Garbage-in-garbage-out applies: the extractor is only as sound
as the min-entropy number it is handed.

\section{The heuristic we started with --- and why it was wrong}
The first implementation used
\[
\Hmin \;\approx\; m\cdot\Fxeb\cdot\big(n-\log_2 n-\log_2\ln 2\big),
\]
giving $\approx1659$ bits at $n=8,m=1000,\Fxeb=0.3$. This was \emph{formulated, not cited}:
three plausible factors multiplied. Only one is legitimate (the per-sample Porter--Thomas
min-entropy $n-\log_2 n-\log_2\ln2$). The other two are wrong:
\begin{itemize}[leftmargin=1.6em,itemsep=1pt]
\item $\times\Fxeb$ treats the score as a per-sample \emph{multiplier} --- but the score is
a \emph{gate} (pass/fail), not a dial.
\item the per-sample constant is the wrong quantity for accumulation (see below).
\end{itemize}

\begin{plain}
The heuristic was an educated placeholder so the pipeline could run --- a formula with the
right variables but invented glue. The danger was never the placeholder; it was forgetting
it was one and quoting its number as if certified. The fix below replaces it with a
published theorem.
\end{plain}

\section{The honest-server bound (our near-term scope)}
\begin{assumption}[Honest server]\label{ass:hs}
Every accepted round is sampled faithfully on the QPU. In the paper's notation this is the
special case $\Qmin=M$.
\end{assumption}

\subsection{Per-round entropy: the exact collision entropy}
\begin{lemma}[Porter--Thomas collision entropy; Liu \emph{et al.}\ SM Eq.~III.33]\label{lem:h2}
\[
H_2(X_i\mid C_i) = -\log_2\!\Big(\int_0^\infty N f(p)\,p^2\,dp\Big)
= -\log_2\!\Big(N^2\tfrac{2}{N^3}\Big) = -\log_2\tfrac{2}{N} = n-1 .
\]
\end{lemma}

\begin{okbox}
The per-round quantity is the \emph{exact} $2$-R\'enyi (collision) entropy, $n-1$ --- not a
conservative von Neumann floor. (An earlier draft of ours justified $n-1$ as a ``von Neumann
floor with $0.39$ bits of slack''; the value was right, the reasoning wrong. The von Neumann
entropy of Porter--Thomas is $n-(1-\gamma)/\ln2\approx n-0.610$ and plays no role here.) The
collision entropy is what the i.i.d.\ min-entropy accumulation requires.
\end{okbox}

\subsection{Accumulation: i.i.d., no EAT, no $\eta$}
Because each quantum round is i.i.d.\ (a stated adversary restriction), the smooth
min-entropy of $Q$ rounds obeys the Tomamichel--Colbeck--Renner i.i.d.\ bound:
\[
\Hmin^{\epss}(X_Q\mid \tilde I_{\mathrm{sn}}) \;\ge\; Q(n-1)-\log_2\tfrac1{\epss}.
\]
There is \emph{no} $\sqrt{M}$ second-order term and \emph{no} generalized-EAT constant
$\eta$ to import --- a seam we once worried about that simply does not appear in the
implemented bound.

\section{The certified bound (Liu \emph{et al.})}
\begin{theorem}[Liu \emph{et al.}\ Theorem~1 / SM Theorem~3, Eq.~III.30]\label{thm:main}
Given $\epss\in(0,\tfrac14)$, the protocol either aborts with probability $>1-4\epss$, or
\[
\boxed{\;\Hmin^{\epss}(X^M\mid\tilde I_{\mathrm{sn}}) \;\ge\; \Qmin\,(n-1)-\log_2\tfrac1{\epss}\;}
\]
\end{theorem}
Extraction with a quantum-proof strong (Toeplitz) extractor gives (SM Corollary~7):
\[
\ell \;\le\; \Qmin(n-1) - 3\log_2\tfrac1{\epssou} - 2 .
\]

\begin{okbox}
\textbf{Verified against the paper.} At $n=56,\ \Qmin=1297,\ \epssou=10^{-6},\ \epss=\epssou/4$:
\[
\Hmin\ge 1297\cdot55-\log_2\tfrac1{2.5\times10^{-7}} = 71{,}335-21.9 = \mathbf{71{,}313}\text{ bits},
\]
matching the paper exactly; and $\ell\le 71{,}335-3\log_2 10^6-2=\mathbf{71{,}273}$ bits,
exactly what they extracted. Our code reproduces both.
\end{okbox}

\paragraph{Our toy honest-server case.} With $\Qmin=M$ (Assumption~\ref{ass:hs}),
$n=8,M=1000,\epss=2^{-33}$:
\[
\Hmin\ge 1000\cdot7-33 = 6967\text{ bits}\quad(\text{vs.\ the heuristic }1659).
\]
The $\approx4\times$ gap is exactly the heuristic's two errors (spurious $\times\Fxeb$, wrong
constant).

\section{$\Qmin$: the real remaining object}
The score $\chi$ and the deadline enter \emph{only} through $\Qmin$, the forced
quantum-round count (Liu SM Lemma~2):
\begin{align*}
\Phi_{\mathcal A}&=\min\!\Big(M-Q,\ \tfrac{c_{\mathrm{eff}}\,\mathcal A\,(M\,t_{\mathrm{th}})}{B}\Big), &
\varepsilon_1&=\exp\!\big(-\tfrac{\delta^2\Phi_{\mathcal A}}{3}\big),\\
L_{\max}&=Q+\Phi_{\mathcal A}(1+\delta), &
\varepsilon_2&=\textstyle\sum_{l=0}^m \mathrm{Hyp}(l;M,L_{\max},m)\,\widetilde\Gamma^{c}(m{+}l,m(\chi{+}1)),
\end{align*}
with $\varepsilon_{\mathrm{adv}}=\min_\delta(\varepsilon_1+\varepsilon_2)$ and
$\Qmin=\min\{Q:\varepsilon_{\mathrm{adv}}\ge4\epss\}$. Here $\mathcal A$ is the adversary's
FLOP rate, $B$ the per-circuit cost, $c_{\mathrm{eff}}$ numerical efficiency.

\begin{plain}
There is no ``score $\to$ entropy curve'' to fill in: the per-round entropy is a flat
$n-1$, and the score's whole role is to set $\Qmin$ --- how many rounds the attacker is
forced to run on real hardware. Computing $\Qmin$ \emph{is} the adversary analysis, and it
is what makes the certified number positive only at large scale.
\end{plain}

\begin{warnbox}
\textbf{$\Qmin$ is transcribed, not yet calibrated.} Our \texttt{compute\_q\_min} implements
the formulas faithfully and is structurally correct (monotone, sensible), but with the
paper's Table-I values it reproduces $\Qmin$ only to order of magnitude ($\sim1700$ vs
$1297$); the exact conventions for $B,c_{\mathrm{eff}},\mathcal A$ and the $\delta$-rule live
in the authors' released code (Zenodo 10.5281/zenodo.12952178). We deliberately did
\emph{not} tune parameters to force $1297$ --- that would be fitting to the answer.
Reconciling against their code turns ``reproduce $1297$'' into a permanent regression test.
\end{warnbox}

\section{Aaronson--Hung vs.\ Liu: which bound we use}
We use the \textbf{Liu \emph{et al.}}\ bound, which is the \emph{corrected, finite-size,
one-sample-per-circuit} descendant of Aaronson--Hung (AH, arXiv:2303.01625). AH is the
asymptotic ancestor: it shows passing XEB forces $\Omega(n)$ bits, but with unknown
finite-size constants, for many-samples-from-one-circuit, and with a technical flaw (its
single-round \emph{min}-entropy bound is impossible --- the $99\%$-honest/$1\%$-zeros
counterexample). Liu corrects it (bounding von Neumann/collision entropy), adapts it to one
sample per circuit, and replaces the asymptotics with a concrete restricted-adversary
analysis that yields actual numbers.

\begin{plain}
Say ``based on the Aaronson--Hung framework,'' but the precise, defensible statement is
``we use Liu \emph{et al.}'s corrected finite-size bound, which adapts Aaronson--Hung to
one-sample-per-circuit.'' AH alone does not give a number you can extract against at finite
scale.
\end{plain}

% =====================================================================
\part*{Part IV --- Status and honest accounting}
\addcontentsline{toc}{section}{Part IV --- Status and honest accounting}
% =====================================================================

\section{What is verified, and what remains}
\begin{center}
\renewcommand{\arraystretch}{1.25}
\begin{tabular}{>{\raggedright}p{0.46\textwidth} p{0.46\textwidth}}
\toprule
\textbf{Verified / closed} & \textbf{Remaining} \\
\midrule
Per-round entropy $H_2=n-1$ (exact, Lemma~\ref{lem:h2}). & Calibrate $\Qmin$ to the authors' code (Mehdi). \\
Certified bound reproduces $71{,}313$ bits. & Real QPU for Leg 4 (Property 1; Colton). \\
Toeplitz extraction reproduces $71{,}273$; extractor is quantum-proof (closes the quantum-conditional-LHL question). & Timing check in Leg 5 (Property 4). \\
No $\eta$/EAT term needed (i.i.d.\ accumulation). & At production scale, $\Qmin$ becomes the central security argument. \\
\bottomrule
\end{tabular}
\end{center}

\section{Honest caveats that did not change}
\begin{warnbox}
\begin{itemize}[leftmargin=1.4em,itemsep=2pt]
\item \textbf{Toy scale is $\approx0$ adversarially.} At $n=8$ a laptop simulates every
circuit, so $\Qmin\to0$ and the certified entropy is $\approx0$ --- by design. The positive
$6967$ exists \emph{only} under the honest-server assumption; quote it with that label.
\item \textbf{The two biggest product gaps are untouched by the math.} Leg 4 is still a
simulator (Property 1) and the timing check is still stubbed (Property 4). Those are
hardware/engineering problems, not math problems.
\item \textbf{Verifiability has limits at scale.} The regime where security switches on is
roughly the regime where brute-force classical verification dies, which is why the protocol
needs the $\Fxeb$ \emph{estimator} and spot-checking rather than full verification.
\end{itemize}
\end{warnbox}

\section{Why security is the whole point}
``Certified quantum'' and ``secure'' are the \emph{same} claim: both reduce to ``no feasible
classical machine could have produced this on demand.'' So the proof-of-quantumness and the
proof-of-security are one proof. Commercially, the only thing worth a premium over a free
classical RNG is the certificate; the load-bearing market question (a Colton question) is
whether customers will pay for \emph{physical} unpredictability over a cheaper
\emph{cryptographic} one (e.g.\ Chainlink VRF). Likely yes in high-stakes, regulated, or
adversarial settings; worth validating with a concrete willingness-to-pay signal before the
expensive QPU spend.

\section{Summary}
\begin{keybox}
The randomness-certification math is now anchored to a published, verified theorem (Liu
\emph{et al.}), reproducing their exact numbers, with two of our own earlier errors corrected
and one phantom ``open seam'' ($\eta$/EAT) dissolved. The single genuinely-open piece is the
$\Qmin$ calibration, now reduced to a bounded regression-test task. This is real progress in
\emph{rigour and trustworthiness} --- not in toy-scale capability, which remains gated on a
real QPU (Leg 4) and the timing check (Leg 5). Every later claim now rests on something
checkable.
\end{keybox}

\vspace{0.5em}
\hrule
\vspace{0.5em}
\small
\noindent\textbf{References.}
M.\ Liu \emph{et al.}, \emph{Certified randomness using a trapped-ion quantum processor},
Nature \textbf{640}:343--348 (2025); arXiv:2503.20498 (Methods + Supplemental Material).\quad
S.\ Aaronson \& S.-H.\ Hung, \emph{Certified randomness from quantum supremacy}, STOC 2023;
arXiv:2303.01625.\quad
M.\ Tomamichel, R.\ Colbeck, R.\ Renner, IEEE Trans.\ Inf.\ Theory \textbf{55}, 5840 (2009).\quad
F.\ Dupuis, O.\ Fawzi, R.\ Renner, Commun.\ Math.\ Phys.\ \textbf{379}, 867 (2020).
\end{document}
